\chapter{Ordres, bons ordres i segments}\label{cap:bons-ordres}

\begin{objectius}
\apartat{Prerequisits} Del capítol~\ref{cap:relacions}: relacions com a parts d'un producte, diagonal, relació inversa. Del capítol~\ref{cap:naturals}: l'ordre de \(\N\), el seu bon ordre, i la convenció que les propietats d'un ordre són afirmacions sobre parts que es comproven sobre punts globals (observació~\ref{obs:ordre-sobre-punts}). Del capítol~\ref{cap:valors}: complements i el lema del buit. Del capítol~\ref{cap:cantor}: la comparació de conjunts, per contrast amb la comparació d'ordres.
\apartat{Objectius} En acabar el capítol, el lector hauria de poder:
\begin{itemize}
  \item definir ordre estricte, ordre total i bon ordre com a parts de \(W\times W\) amb condicions, i saber quines es llegeixen sobre punts;
  \item demostrar que tot bon ordre és total i que el mínim és únic;
  \item definir els segments inicials \(W_{<w}\) i les parts inicials, i demostrar que en un bon ordre les parts inicials pròpies són exactament els segments, que són ben ordenats i determinen el seu punt;
  \item demostrar la rigidesa (cap punt baixa sota una creixent injectiva), la unicitat de l'isomorfisme d'ordre i que cap bon ordre és isomorf a un segment propi;
  \item definir la comparació \(\preceq_{\mathrm o}\) per segments, demostrar-ne l'antisimetria i la transitivitat, i la part d'exclusivitat de la tricotomia, sabent que la comparabilitat total demana recursió transfinita i no elecció.
\end{itemize}
\end{objectius}

Amb el llenguatge de les famílies a punt (capítol~\ref{cap:families}), la
part transfinita necessita ara la noció que en serà el fil conductor: la
d'\emph{ordre}, i sobretot la de \emph{bon ordre}. Un bon ordre és la forma
abstracta d'un procés que es recorre \enquote{de baix cap amunt} sense
saltar-se cap pas ni quedar-se encallat: tota subcol·lecció no buida té un
primer element. Els naturals ja en van ser el primer exemple
(teorema~\ref{teo:bon-ordre-N}); el que fem aquí és desprendre la propietat
de \(\N\) i estudiar-la sobre un conjunt qualsevol. La recompensa és doble.
D'una banda, els bons ordres tenen una \emph{rigidesa} notable ---entre dos
n'hi ha, com a molt, un isomorfisme d'ordre--- que els fa servir de
\emph{mesura de forma}: dir que un és \enquote{més curt} que un altre serà,
literalment, que és isomorf a un segment inicial de l'altre. De l'altra,
aquesta rigidesa és exactament el que permetrà, al capítol següent, definir
funcions per recursió al llarg d'un bon ordre. Aquest capítol estableix les
definicions ---sempre primer sobre parts, després sobre punts--- i els cinc
resultats bàsics sobre segments i comparació, sense cap axioma nou.

\section{Ordres estrictes, totals i bons ordres}\label{sec:ordres}

Com totes les relacions del llibre (capítol~\ref{cap:relacions}), un ordre
no serà una fórmula sinó una \emph{part} d'un producte; les lectures sobre
punts en són la pertinença generalitzada. Mantenim la convenció del
capítol~\ref{cap:naturals} (observació~\ref{obs:ordre-sobre-punts}): les
propietats d'un ordre són afirmacions sobre parts de \(W\times W\) o de
\(W\), i per la separació global i el lema del buit n'hi ha prou amb
comprovar-les sobre punts globals, on \enquote{o} i \enquote{existeix}
tenen el sentit clàssic. Les condicions que es formulen amb una implicació
(com la transitivitat) són inclusions de parts i valen a tots els estadis;
les que es formulen amb una disjunció (com la totalitat) són igualtats de
parts amb una unió, i sobre un estadi general un element pot quedar
repartit entre els membres.

\begin{definicio}[Ordre estricte]\index{Ordre estricte}\label{def:ordre-estricte}
Un \emph{ordre estricte} sobre un conjunt \(W\) és una part
\[
{<}\ \hookrightarrow\ W\times W
\]
que és \emph{irreflexiva} i \emph{transitiva}. Escrivim \(u<w\) per a la
pertinença generalitzada \((u,w)\mathrel{\varepsilon}{<}\)
(capítol~\ref{cap:relacions}). Les dues condicions, sobre elements
generalitzats \(u,v,w\colon X\to W\) a un estadi qualsevol, són:
\begin{itemize}
\item[(o1)] \emph{irreflexivitat}: mai \(w<w\); és a dir, la part \({<}\) és
disjunta de la diagonal \(\Delta\hookrightarrow W\times W\)
(capítol~\ref{cap:relacions});
\item[(o2)] \emph{transitivitat}: si \(u<v\) i \(v<w\), aleshores \(u<w\).
\end{itemize}
El parell \((W,<)\) és un \emph{conjunt ordenat}.
\end{definicio}

\begin{observacio}[L'ordre és una part, no un esquema]\label{obs:ordre-part-general}
Com per a \(\N\) (observació~\ref{obs:ordre-part}), \({<}\) és una part
concreta de \(W\times W\); \enquote{\(u<w\)} n'és la pertinença
generalitzada, i les condicions (o1)--(o2) es tallen amb la diagonal, la
composició de relacions i els quantificadors interns
(capítols~\ref{cap:relacions} i~\ref{cap:parts-interna}). De \({<}\) se'n
deriva la relació \emph{no estricta} \(\le\), la part
\({\le}:={<}\cup\Delta\hookrightarrow W\times W\), amb \(u\le w\Leftrightarrow
(u<w\ \vee\ u=w)\); és reflexiva, i és transitiva i antisimètrica quan
\({<}\) és un ordre estricte (la irreflexivitat de \({<}\) dona
l'antisimetria de \(\le\)). Treballarem indistintament amb \({<}\) i amb
\(\le\), segons convingui.
\end{observacio}

\begin{definicio}[Ordre total]\index{Ordre total}\label{def:ordre-total-gen}
Un ordre estricte \({<}\) sobre \(W\) és \emph{total} (o \emph{lineal}) si
la unió \(\Delta\cup{<}\cup\operatorname{Inv}(<)\) és la part total de
\(W\times W\); sobre punts globals, si dos punts qualssevol són
comparables:
\[
\text{per a tots } u,w\colon\mathbf1\to W,\qquad u=w\ \ \text{o}\ \ u<w\ \ \text{o}\ \ w<u,
\]
i aquestes tres possibilitats són mútuament excloents (\emph{tricotomia}
puntual), la darrera exclusió per la irreflexivitat i la transitivitat.
\end{definicio}

\begin{definicio}[Bon ordre]\index{Bon ordre}\label{def:bon-ordre}
Un ordre estricte \({<}\) sobre \(W\) és un \emph{bon ordre}, i \((W,<)\) un
conjunt \emph{ben ordenat}, si tota part no inicial de \(W\) té un element
mínim: per a tota part \(P\hookrightarrow W\) no inicial existeix un punt
\(m\colon\mathbf1\to W\) que es factoritza per \(P\) i compleix \(m\le k\)
per a tot element generalitzat \(k\) de \(P\). El mínim, quan existeix, és
únic.
\end{definicio}

\begin{lema}[Les parts es determinen pels punts]\label{lema:transferencia-punts}
Siguin \(R,S\hookrightarrow W\times W\) dues parts (per exemple, dues relacions
sobre \(W\)). Sota la separació global (teoria ben apuntada,
axioma~\ref{ax:separacio-globals}), si \(R\) i \(S\) tenen exactament els
mateixos \emph{punts globals} ---per a tot \(p\colon\mathbf1\to W\times W\),
\(p\) es factoritza per \(R\) si i només si es factoritza per \(S\)---,
aleshores \(R\equiv S\) com a parts, i per tant coincideixen sobre \emph{tot}
element generalitzat \(X\to W\times W\), sigui quin sigui l'estadi \(X\). En
particular, una relació que conté tots els punts globals de \(W\times W\)
és la part total.
\end{lema}

\begin{prova}
És la reducció als punts que ja vam fer al capítol~\ref{cap:relacions}
(pràctica) i al capítol~\ref{cap:naturals}: la part \(R\cap S^c\) no té cap
punt i, pel lema del buit (lema~\ref{lema:buit}), és buida, de manera que
\(R\subseteq S\) per la descomposició \(R\cong(R\cap S)+(R\cap S^c)\); i
simètricament. Dues parts incloses l'una en l'altra són equivalents, i la
pertinença d'un element generalitzat es decideix per factorització, igual
en totes dues. Observeu el que el lema \emph{no} diu: que una disjunció
verificada sobre punts (\enquote{\(u<w\) o \(w<u\)}) valgui, disjunt per
disjunt, sobre elements generalitzats; el que val és que la \emph{unió} de
les parts corresponents és total.
\end{prova}

\begin{proposicio}[Tot bon ordre és total]\label{prop:bon-ordre-total}
Si \({<}\) és un bon ordre sobre \(W\), aleshores és total. El mínim d'una
part no inicial és, a més, únic.
\end{proposicio}

\begin{prova}
La unicitat del mínim surt de l'antisimetria de \(\le\)
(observació~\ref{obs:ordre-part-general}): dos mínims \(m,m'\) compleixen
\(m\le m'\) i \(m'\le m\), d'on \(m=m'\). Per a la totalitat, provem primer la
comparabilitat sobre \emph{punts globals} i després la transferim a tots els
estadis. Siguin \(u,w\colon\mathbf1\to W\) dos punts distints i considerem la
part \(\{u,w\}\hookrightarrow W\) ---la imatge del parell
\(\langle u,w\rangle\colon\mathbf1+\mathbf1\to W\)
(teorema~\ref{teo:factoritzacio})---, que és no inicial. Pel bon ordre, té un
mínim, que és \(u\) o bé \(w\). Si el mínim és \(u\), aleshores \(u\le w\), i
com que \(u\neq w\), \(u<w\); si és \(w\), simètricament \(w<u\). Així la part
\(\Delta\cup{<}\cup\operatorname{Inv}(<)\hookrightarrow W\times W\) conté
tots els punts globals de \(W\times W\). Pel
lema~\ref{lema:transferencia-punts} (separació global), aquesta part és tot
\(W\times W\), que és la totalitat de la definició~\ref{def:ordre-total-gen}.
\end{prova}

\begin{observacio}[Lectura clàssica]\label{obs:bon-ordre-classic}
Un bon ordre és un ordre total en què tot subconjunt no buit té un primer
element. La condició, aparentment modesta, és molt forta: prohibeix les
cadenes descendents infinites \(w_0>w_1>w_2>\cdots\) ---una tal cadena tindria
per imatge una part sense mínim--- i és el que fa possible tant la inducció
com la recursió al llarg de \(W\). Els naturals amb el seu ordre habitual en
són l'exemple canònic (teorema~\ref{teo:bon-ordre-N}); en canvi \(\Z\), \(\Q\)
i \(\R\) amb l'ordre usual \emph{no} són ben ordenats, perquè tenen parts no
buides sense mínim (el mateix \(\Z\), o l'interval obert). Que \emph{tot}
conjunt admeti \emph{algun} bon ordre és una afirmació molt més forta ---el
teorema del bon ordre--- que requereix l'elecció i que veurem al
capítol~\ref{cap:zorn}.
\end{observacio}

\section{Segments inicials}\label{sec:segments}

La noció que organitza tota la comparació de bons ordres és la de segment
inicial: el tros de \(W\) que queda \enquote{estrictament per sota} d'un
punt.

\begin{definicio}[Segment inicial]\index{Segment inicial}\label{def:segment}
Sigui \((W,<)\) un conjunt ordenat i \(w\colon\mathbf1\to W\) un punt. El
\emph{segment inicial} determinat per \(w\) és la part
\[
W_{<w}\ :=\ \{\,u\in W\ \mid\ u<w\,\}\ \hookrightarrow\ W,
\]
la imatge inversa de la relació \({<}\) al llarg de
\(\langle\id_W,w\circ{!}_W\rangle\); en la lectura sobre punts, els elements
de \(W_{<w}\) són els \(u\) amb \(u<w\). Més generalment, una part
\(S\hookrightarrow W\) és un \emph{segment inicial} (o \emph{part inicial})
de \(W\) si és tancada cap avall: si \(u<w\) i \(w\in S\), aleshores \(u\in
S\).
\end{definicio}

\begin{observacio}[Els \(W_{<w}\) i les parts inicials]\label{obs:segments-inicials}
Tot segment \(W_{<w}\) és una part inicial: si \(u<w'\) i \(w'<w\), la
transitivitat dona \(u<w\). El recíproc, per a bons ordres, és gairebé cert i
molt útil: veurem (proposició~\ref{prop:parts-inicials}) que en un bon ordre
les úniques parts inicials pròpies són exactament els segments \(W_{<w}\).
Aquesta identificació ---parts inicials \(=\) segments--- és la que farà que
\enquote{ser un segment de} funcioni com una relació d'ordre entre bons
ordres.
\end{observacio}

\begin{proposicio}[Els segments d'un bon ordre són ben ordenats]\label{prop:segment-ben-ordenat}
Sigui \((W,<)\) un bon ordre i \(S\hookrightarrow W\) una part qualsevol.
Aleshores la restricció de \({<}\) a \(S\) ---la part
\({<}\cap(S\times S)\hookrightarrow S\times S\)--- és un bon ordre sobre
\(S\). En particular, tot segment inicial \(W_{<w}\) és ben ordenat.\par
\emph{Axiomes usats:} parts i interseccions (capítols~\ref{cap:parts}
i~\ref{cap:valors}); cap axioma nou.
\end{proposicio}

\begin{prova}
La restricció \({<}\!\restriction_S\,={<}\cap(S\times S)\) és una part de
\(S\times S\) (intersecció de parts, capítol~\ref{cap:valors}), i és
irreflexiva i transitiva perquè ho és \({<}\): les condicions (o1)--(o2) es
conserven en restringir a una part. Per al bon ordre, sigui
\(P\hookrightarrow S\) una part no inicial; via la inclusió
\(S\hookrightarrow W\), \(P\) és també una part no inicial de \(W\)
(la composició de monomorfismes és un monomorfisme,
proposició~\ref{prop:preceq-preordre}), que pel bon ordre de \(W\) té un
mínim \(m\). Com que \(m\in P\subseteq S\) i \(m\le k\) per a tot \(k\in
P\), aquest \(m\) és també el mínim de \(P\) dins \(S\). Per tant
\({<}\!\restriction_S\) és un bon ordre. El cas \(S=W_{<w}\) n'és una
instància.
\end{prova}

\begin{proposicio}[Les parts inicials pròpies són segments]\label{prop:parts-inicials}
Sigui \((W,<)\) un bon ordre i \(S\subsetneq W\) una part inicial pròpia.
Aleshores existeix un únic punt \(w\colon\mathbf1\to W\) tal que
\(S=W_{<w}\); concretament, \(w\) és el mínim del complement
\(W\setminus S\).
\end{proposicio}

\begin{prova}
Com que \(S\) és pròpia, el complement \(W\setminus S\)
(teorema~\ref{teo:part-complementari}) és no inicial, i pel bon ordre té un
mínim \(w\). Vegem \(S=W_{<w}\). Si \(u\in S\), aleshores \(u<w\): altrament
\(w\le u\), i com que \(S\) és inicial i \(u\in S\), de \(w\le u\) en sortiria
\(w\in S\) ---si \(w<u\) per inicialitat, i si \(w=u\) trivialment---, contra
\(w\in W\setminus S\); així \(u<w\), és a dir \(u\in W_{<w}\).
Recíprocament, si \(u<w\), aleshores \(u\in S\): com que \(w\) és el
\emph{mínim} de \(W\setminus S\), cap element estrictament menor que \(w\)
és al complement, de manera que \(u\in S\). Per tant \(S=W_{<w}\). La
unicitat de \(w\) és la del punt que determina un segment donat, que provem
tot seguit (proposició~\ref{prop:segment-determina-punt}).
\end{prova}

\begin{proposicio}[Un segment determina el seu punt]\label{prop:segment-determina-punt}
Sigui \((W,<)\) un bon ordre. Si \(W_{<v}=W_{<w}\) com a parts de \(W\),
aleshores \(v=w\). És a dir, dos segments inicials iguals ---i, més avall,
isomorfs--- corresponen al mateix punt.
\end{proposicio}

\begin{prova}
Suposem \(W_{<v}=W_{<w}\) i, cap a una contradicció, \(v\neq w\). Per la
totalitat (proposició~\ref{prop:bon-ordre-total}), \(v<w\) o \(w<v\); diguem
\(v<w\). Aleshores \(v\in W_{<w}=W_{<v}\), és a dir \(v<v\), contra la
irreflexivitat. Simètricament si \(w<v\). Per tant \(v=w\).
\end{prova}

\section{Aplicacions creixents i isomorfismes d'ordre}\label{sec:creixents}

\begin{definicio}[Aplicació creixent i isomorfisme d'ordre]\index{Aplicació creixent i isomorfisme d'ordre}\label{def:creixent}
Siguin \((W,<_W)\) i \((V,<_V)\) conjunts ordenats. Una aplicació
\(f\colon W\to V\) és \emph{creixent} (o \emph{monòtona}) si conserva
l'ordre estricte: \(u<_W w\) implica \(f(u)<_V f(w)\), per a tots els
elements generalitzats \(u,w\colon X\to W\); com a parts,
\({<_W}\subseteq(f\times f)^{-1}({<_V})\), una inclusió, que val a tots
els estadis si val sobre punts. És un \emph{isomorfisme
d'ordre} si és un isomorfisme de conjunts (definició~\ref{def:isomorfisme})
i la seva inversa també és creixent; equivalentment, si és una bijecció amb
\(u<_W w\Leftrightarrow f(u)<_V f(w)\). Escrivim \((W,<_W)\cong(V,<_V)\) quan
n'existeix un.
\end{definicio}

\begin{observacio}[Creixent i injectiva]\label{obs:creixent-injectiva}
Sobre un ordre total ---en particular un bon ordre--- tota aplicació
creixent és automàticament un monomorfisme: sobre punts, si \(f(u)=f(w)\)
amb \(u\neq w\), la totalitat dona \(u<w\) o \(w<u\), i la monotonia
\(f(u)<f(w)\) o \(f(w)<f(u)\), contra \(f(u)=f(w)\) i la irreflexivitat;
injectiva sobre punts és mono (observació~\ref{obs:injectiva-mono}). Per tant, entre
conjunts totalment ordenats, un isomorfisme d'ordre no és més que una
aplicació creixent, exhaustiva i amb inversa creixent; i, com veurem, entre
bons ordres n'hi ha com a molt un.
\end{observacio}

Els isomorfismes d'ordre entre bons ordres són \emph{rígids}: no en poden
existir dos diferents. Aquest és el resultat que fa dels bons ordres una
bona mesura de forma, i el que garantirà, al capítol següent, la unicitat de
les definicions per recursió transfinita.

\begin{lema}[Rigidesa: cap punt es mou cap avall]\label{lema:rigidesa}
Sigui \((W,<)\) un bon ordre i \(f\colon W\to W\) una aplicació creixent i
injectiva. Aleshores \(w\le f(w)\) per a tot \(w\colon\mathbf1\to W\); cap
element no baixa.
\end{lema}

\begin{prova}
Sigui \(P:=\{\,w\in W\mid f(w)<w\,\}\hookrightarrow W\) la part dels punts
que \emph{baixen} ---la imatge inversa per \(\langle f,\id_W\rangle\) de la
relació \({<}\) (capítols~\ref{cap:parts} i~\ref{cap:relacions}). Volem
\(P\) inicial. Si no ho fos, pel bon ordre tindria un mínim \(w_0\), amb
\(f(w_0)<w_0\). Aplicant \(f\), que és creixent, a la desigualtat
\(f(w_0)<w_0\) obtenim \(f(f(w_0))<f(w_0)\); és a dir, el punt
\(w_1:=f(w_0)\) també baixa, \(w_1\in P\), i \(w_1=f(w_0)<w_0\) és
estrictament menor que el mínim \(w_0\) de \(P\), una contradicció. Per tant
\(P\) és inicial: cap punt baixa, i \(w\le f(w)\) per a tot \(w\).
\end{prova}

\begin{teorema}[Unicitat de l'isomorfisme d'ordre]\label{teo:unicitat-iso}
Entre dos bons ordres \((W,<_W)\) i \((V,<_V)\) hi ha, com a molt, un
isomorfisme d'ordre. En particular, l'únic automorfisme d'ordre d'un bon
ordre és la identitat.\par
\emph{Axiomes usats:} bon ordre, parts, imatges inverses
(capítols~\ref{cap:parts} i~\ref{cap:imatges}); cap axioma nou; cap elecció.
\end{teorema}

\begin{prova}
Siguin \(f,g\colon W\to V\) dos isomorfismes d'ordre. Aleshores
\(g^{-1}\circ f\colon W\to W\) és un automorfisme d'ordre (composició
d'isomorfismes d'ordre), i també ho és \(f^{-1}\circ g\), la seva inversa.
Pel lema~\ref{lema:rigidesa} aplicat a \(g^{-1}\circ f\), \(w\le
(g^{-1}\circ f)(w)\) per a tot \(w\); aplicat a \(f^{-1}\circ g\),
\(w\le(f^{-1}\circ g)(w)\), i component amb \(g^{-1}\circ f\)
(creixent) s'obté \((g^{-1}\circ f)(w)\le w\). Les dues desigualtats donen
\((g^{-1}\circ f)(w)=w\) per a tot \(w\), és a dir \(g^{-1}\circ f=\id_W\),
d'on \(f=g\). El cas \(W=V\) dona que l'únic automorfisme és la identitat.
\end{prova}

\begin{corollari}[Un bon ordre no és isomorf a cap segment propi]\label{cor:no-segment-propi}
Cap bon ordre \((W,<)\) no és isomorf a cap dels seus segments inicials
propis \(W_{<w}\). Més en general, cap isomorfisme d'ordre no aplica \(W\)
sobre una part inicial pròpia.
\end{corollari}

\begin{prova}
Suposem un isomorfisme d'ordre \(f\colon W\cong W_{<w}\) per a algun \(w\).
Compost amb la inclusió \(W_{<w}\hookrightarrow W\), dona una aplicació
creixent i injectiva \(f\colon W\to W\) amb \(f(w)\in W_{<w}\), és a dir
\(f(w)<w\). Però això contradiu el lema~\ref{lema:rigidesa}, segons el qual
\(w\le f(w)\). Per tant no existeix cap tal isomorfisme. El cas d'una part
inicial pròpia qualsevol es redueix a un segment
(proposició~\ref{prop:parts-inicials}).
\end{prova}

\section{Comparació de bons ordres}\label{sec:comparacio-bons-ordres}

Els resultats anteriors permeten \emph{comparar} bons ordres per la relació
\enquote{ser isomorf a un segment inicial de}. Aquesta relació es comporta
com un ordre, i és la que, al capítol dels tipus ordinals
(capítol~\ref{cap:ordinals}), es convertirà en la relació \(<\) entre ordinals.

\begin{definicio}[Un bon ordre com a segment d'un altre]\index{Un bon ordre com a segment d'un altre}\label{def:mes-curt}
Direm que \((W,<_W)\) és \emph{(isomorf a) un segment inicial} de
\((V,<_V)\), i escriurem \(W\preceq_{\mathrm o} V\), si existeix un
isomorfisme d'ordre de \(W\) sobre un segment inicial de \(V\) ---sobre tot
\(V\), o sobre algun \(V_{<v}\). Escrivim \(W\prec_{\mathrm o}V\) quan és
sobre un segment \emph{propi} \(V_{<v}\), i \(W\cong V\) quan és sobre tot
\(V\).
\end{definicio}

\begin{proposicio}[La comparació per segments és antisimètrica i transitiva]\label{prop:comparacio-ordres}
Sobre bons ordres:
\begin{enumerate}
\item si \(W\preceq_{\mathrm o}V\) i \(V\preceq_{\mathrm o}W\), aleshores
\(W\cong V\);
\item si \(W\preceq_{\mathrm o}V\) i \(V\preceq_{\mathrm o}U\), aleshores
\(W\preceq_{\mathrm o}U\);
\item l'isomorfisme d'ordre de qualsevol de les situacions és \emph{únic}
(teorema~\ref{teo:unicitat-iso}).
\end{enumerate}
\emph{Axiomes usats:} els dels resultats anteriors; cap axioma nou; cap
elecció.
\end{proposicio}

\begin{prova}
(2) Si \(f\colon W\cong V_{\le}\) sobre un segment de \(V\) i \(g\colon
V\cong U_{\le}\) sobre un segment de \(U\), aleshores \(g\) aplica el segment
\(V_{\le}\) sobre un segment de \(U\). En efecte, sigui \(T\) la part
inicial de \(U\) sobre la qual \(g\) és un isomorfisme d'ordre, i \(S\) una
part inicial de \(V\). Si \(u<_U g(s)\) amb \(s\in S\), aleshores \(u\in T\)
perquè \(T\) és inicial, de manera que \(u=g(v)\) per a un \(v\in V\); com que
\(g\) també \emph{reflecteix} l'ordre (la seva inversa és creixent),
\(v<_V s\), i \(v\in S\) perquè \(S\) és inicial. Així \(g(S)\) és una part
inicial de \(U\) (proposició~\ref{prop:parts-inicials}). Per tant \(g\circ f\) aplica
\(W\) isomòrficament sobre un segment de \(U\); així \(W\preceq_{\mathrm
o}U\). (1) Si \(W\preceq_{\mathrm o}V\preceq_{\mathrm o}W\), per (2)
\(W\preceq_{\mathrm o}W\) via un isomorfisme d'ordre de \(W\) sobre un
segment inicial de \(W\); pel corol·lari~\ref{cor:no-segment-propi} aquest
segment no pot ser propi, de manera que és tot \(W\) i la composició és un
automorfisme ---la identitat (teorema~\ref{teo:unicitat-iso}). Desfent-la,
cadascun dels dos isomorfismes és sobre tot l'altre, d'on \(W\cong V\). (3)
és el teorema~\ref{teo:unicitat-iso}.
\end{prova}

Falta el resultat central de comparació: que dos bons ordres qualssevol
\emph{sempre} són comparables ---un dels dos és isomorf a un segment de
l'altre. És la \emph{tricotomia}, i és la peça que, en el llenguatge
estructural, substitueix moltes manipulacions clàssiques amb ordinals de von
Neumann. La seva prova, però, construeix l'isomorfisme de comparació
\emph{punt per punt al llarg del bon ordre}, i això és una recursió
transfinita, que encara no tenim. Enunciem la tricotomia aquí, amb la seva
dependència declarada, i en donem la part que \emph{no} depèn de la recursió;
la construcció de l'isomorfisme es completa al capítol següent.

\begin{teorema}[Tricotomia de la comparació; part sense recursió]\label{teo:tricotomia}
Siguin \((W,<_W)\) i \((V,<_V)\) dos bons ordres. Aleshores \emph{com a molt
una} de les tres situacions es dona:
\[
W\prec_{\mathrm o}V,
\qquad
W\cong V,
\qquad
V\prec_{\mathrm o}W,
\]
i quan es dona, l'isomorfisme d'ordre corresponent és únic. Que \emph{almenys
una} es dona sempre ---la comparabilitat total--- es demostra construint un
isomorfisme d'ordre entre \(W\) i un segment de \(V\), o a l'inrevés, per
recursió transfinita; ho establirem al teorema~\ref{teo:comparabilitat} del
capítol~\ref{cap:recursio}.\par
\emph{Axiomes usats (per a l'exclusivitat):} unicitat de l'isomorfisme
d'ordre i no-isomorfisme amb segments propis
(teorema~\ref{teo:unicitat-iso}, corol·lari~\ref{cor:no-segment-propi}); cap
elecció. La comparabilitat total usa la recursió
transfinita del capítol~\ref{cap:recursio}, també sense elecció.
\end{teorema}

\begin{prova}[Exclusivitat mútua]
Provem que les tres situacions són incompatibles dues a dues; l'existència
d'almenys una es remet al capítol~\ref{cap:recursio}.
Suposem \(W\cong V\) i
alhora \(W\prec_{\mathrm o}V\); aleshores, component, \(V\cong W\prec_{\mathrm
o}V\) dona \(V\) isomorf a un segment propi de si mateix, contra el
corol·lari~\ref{cor:no-segment-propi}. Igual amb \(W\cong V\) i
\(V\prec_{\mathrm o}W\). Finalment, si \(W\prec_{\mathrm o}V\) i
\(V\prec_{\mathrm o}W\) alhora, la transitivitat
(proposició~\ref{prop:comparacio-ordres}) dona \(W\prec_{\mathrm o}W\), un
isomorfisme de \(W\) sobre un segment propi, altre cop contra el
corol·lari~\ref{cor:no-segment-propi}. Per tant les tres són mútuament
excloents, i la unicitat de l'isomorfisme és el
teorema~\ref{teo:unicitat-iso}.
\end{prova}

\begin{observacio}[Per què la comparabilitat necessita recursió, i per què no elecció]\label{obs:tricotomia-recursio}
La direcció que aquí deixem pendent ---que almenys una de les tres situacions
es dona--- no és una mancança de la nostra prova, sinó un fet estructural:
per decidir quin dels dos bons ordres és \enquote{més curt} cal
\emph{aparellar-los} element a element començant pels mínims respectius, i
aquest aparellament es defineix per recursió al llarg d'un dels dos ordres.
No hi intervé cap \emph{tria}: a cada pas, l'element que toca és el
\emph{mínim} del que queda, que és únic i canònic (com als subobjectes de
\(\N\), observació~\ref{obs:minim-vs-eleccio}); per això la comparabilitat de
bons ordres, a diferència de la comparabilitat de conjunts \emph{qualssevol}
(que veurem amb el teorema del bon ordre), \emph{no} usa l'elecció. El que necessita és
únicament el principi de recursió transfinita amb codomini fix, que
construïm al capítol~\ref{cap:recursio}; per això n'hem separat l'enunciat.
Aquesta és, també, la disciplina que la part transfinita mantindrà: enunciar
un resultat només quan les seves eines ja existeixen, o declarar
explícitament de quina depèn.
\end{observacio}

\begin{resumcap}
\apartat{Idees centrals}
\begin{enumerate}
  \item Un ordre estricte és una part \({<}\hookrightarrow W\times W\) irreflexiva i transitiva; total quan \(\Delta\cup{<}\cup\operatorname{Inv}(<)\) és tot \(W\times W\); bon ordre quan tota part no inicial té mínim. Les propietats es comproven sobre punts, i les disjunctives s'entenen com a unions.
  \item Tot bon ordre és total i el mínim és únic; els segments \(W_{<w}\) són les parts inicials pròpies, són ben ordenats i determinen el seu punt.
  \item Rigidesa: sota una creixent injectiva d'un bon ordre en si mateix cap punt baixa. D'aquí, entre dos bons ordres hi ha com a molt un isomorfisme d'ordre, l'únic automorfisme és la identitat, i cap bon ordre és isomorf a un segment propi.
  \item \(W\preceq_{\mathrm o}V\) (\(W\) isomorf a un segment inicial de \(V\)) és antisimètrica llevat d'isomorfisme i transitiva, i les tres situacions \(W\prec_{\mathrm o}V\), \(W\cong V\), \(V\prec_{\mathrm o}W\) s'exclouen. Que almenys una es doni demana recursió transfinita, no elecció.
\end{enumerate}
\apartat{Errors típics} Llegir la totalitat com \enquote{\(u<w\) o \(w<u\)} per a elements generalitzats (val sobre punts; com a parts és una unió total); creure que un bon ordre pot tenir automorfismes no trivials (com els ordres totals densos); confondre \(\preceq_{\mathrm o}\) (segments, forma) amb \(\preceq\) (monomorfismes, mida) del capítol~\ref{cap:cantor}; i pensar que la comparabilitat de bons ordres usa l'elecció, quan el que usa és el mínim.
\apartat{Lectura recomanada} \cite[§§14 i 17]{halmos} per als conjunts ordenats, els bons ordres, els segments i la comparabilitat en la teoria clàssica, que és el que aquí es reconstrueix sobre parts; \cite[cap.~2]{jech2003} per a la teoria clàssica dels bons ordres i dels ordinals; \cite[cap.~I.6--I.7]{kunen2011} per a la presentació axiomàtica estàndard dels bons ordres i la recursió; i \cite[§§8.4--8.5]{lawvere-rosebrugh} per a l'ordre i els bons ordres en conjunts abstractes.
\end{resumcap}
